\section{Formal Verification}\label{sec:formal-verification}

Shield Reduction (\Cref{thm:shield-reduction}) is formally verified in Lean
4 via Aristotle with no remaining sorries.  For the exact
$5n/24+O(1)$ cover theorem (\Cref{thm:524-cover}), the explicit cover/packing
identities are likewise formalized with no remaining sorries, while the final
floor-sum cardinality count remains elementary prose.  The $13/36$ upper bound
(\Cref{thm:1336}) is likewise formally verified with no remaining sorries.
Theorem~A (\Cref{thm:theorem-a}) and the $5/16$ upper bound
(\Cref{thm:516}) have Lean artifacts with classical-number-theory sorries
(PNT- and Chebyshev-$\vartheta$-adjacent, enumerated in
\Cref{app:formalization}); these are Mathlib-dependency gaps rather than
logical holes.  T2 has a Lean-verified finite core with the activation-stage
wrapper pending.  The endgame reduction
for the Bonferroni--4 theorem --- namely, that if the four Round~15 moment
terms are componentwise within $10^{-5}$ of their certified limits
$J_1,\dots,J_4$ plus eventual $o(1)$ error, then $\L(n)<0.19n$ eventually --- is
formally verified in Lean~4 with no remaining sorries in
\path{erdos_872_core/RequestProject/Round15Bonferroni4/Target.lean}.  The
supplying moment-convergence hypotheses of \Cref{thm:bonf-comp,thm:prime-bridge}
are proved rigorously in prose in \Cref{sec:main-upper}, with multi-verifier
audit and sandbox computation.

All paths in \Cref{tab:lean-summary} are relative to the public artifact
repository \cite{ErdosHarness26}, with the Lean-specific paths rooted at
\texttt{erdos-872/lean/}.  \Cref{app:formalization} lists the current explicit
\texttt{sorry} sites in the checked-in files rather than the older aggregate
counts recorded in the repository index.

\begin{table}[ht]
\centering
\footnotesize
\renewcommand{\arraystretch}{1.15}
\begin{tabular}{|p{0.20\textwidth}|p{0.33\textwidth}|p{0.09\textwidth}|p{0.27\textwidth}|}
\hline
\textbf{Result} & \textbf{Lean artifact} & \textbf{Sorries} & \textbf{Status note} \\
\hline
Shield Reduction &
\path{shield_reduction/shield_reduction_aristotle/RequestProject/ShieldReduction.lean} &
0 &
Complete formal proof of the terminal inequality \\
\hline
Theorem A / T1 &
\path{theorem_A/theorem_A_shield_lower_bound_aristotle/RequestProject/ShieldMainTheorem.lean} &
5 &
All remaining sorries are PNT-adjacent analytic inputs \\
\hline
Exact $5/24$ cover &
\path{tau_5_24/tau_5_24_aristotle/RequestProject/\{Cover,Packing,Tau\}.lean} &
0 &
Structural identities are formalized; the floor-sum cardinality count remains elementary prose \\
\hline
$13/36$ upper bound &
\path{shortener_13_36/shortener_13_36_v2_aristotle/RequestProject/Shortener/MainTheorem.lean} &
0 &
Complete formal proof for the truncation-based strategy \\
\hline
$5/16$ upper bound &
\path{shortener_5_16/shortener_5_16_aristotle/RequestProject/Shortener516/Theorems.lean} &
1 &
One bundled Chebyshev / prefix-legality gap remains; the combinatorial core is zero-sorry \\
\hline
T2 finite core &
\path{erdos_872_core/RequestProject/T2Finite/*.lean} &
0 in current files &
Graph/hypergraph potentials and residual comparison are formalized; the activation-stage wrapper is still being integrated \\
\hline
Round~15 endgame reduction &
\path{erdos_872_core/RequestProject/Round15Bonferroni4/Target.lean} &
0 &
Formalizes \texttt{eventually\_strict\_lt\_point19\_of\_componentwise\_close}; the analytic moment convergence is supplied by \Cref{thm:bonf-comp,thm:prime-bridge} \\
\hline
\end{tabular}
\caption{Current Lean / Aristotle verification status.}
\label{tab:lean-summary}
\end{table}

\begin{remark}
Two distinctions are important.  First, a zero-sorry reduction is not the same
thing as a full end-to-end formalization of the surrounding theorem: the
Round~15 Lean file closes the last step of the proof, while the analytic
moment-convergence input remains in prose.  Second, the T2 project already
contains zero-sorry formalizations of the finite potential game and residual
comparison layers, but not yet a single final Lean theorem packaging the full
activation-stage argument of \Cref{thm:t2}.
\end{remark}
